% Part:sets-functions-relations
% Chapter: size-of-sets
% Section: equinumerous-sets

\documentclass[../../../include/open-logic-section]{subfiles}

\begin{document}

\olfileid{sfr}{siz}{equ}

\olsection{همشمېر والی}

موږ د سټونو د «اندازې» يوه شهودي مفکوره لرو چې د متناهي سټونو
دپاره سمه کار کوي۔ خو د نامتناهي سټونو څه؟ که وغواړو د \emph{هرې}
اندازې د دوو سټونو د اندازو د پرتله کولو يوه صوري لار جوړه کړو،
ښه به وي چې لومړے تعريف کړو سټونه کله يوه اندازه لري۔ دا ئې د
فرېګه وينا ده:
\begin{quote}
  که يو خدمتګار غواړي ډاډه شي چې پر مېز ئې د پليټونو هومره چاړې
  په دقيقه توګه اېښې دي، د دواړو له شمېرلو پرته هم دا کار کولے شي:
  يوازې د هر پليټ ښي لوري ته يوه چاړه کېږدي، داسې چې پر مېز هره
  چاړه د کوم پليټ ښي لوري ته پرته وي۔ په دې توګه پليټونه او چاړې
  په ځانګړې توګه يو له بل سره، او په رښتيا هم د هماغې مکاني اړيکې
  له لارې، تړل کېږي۔ \citep[\S70]{Frege1884}
\end{quote}
د دې عبارت بصيرت په يوه صوري تعريف کښې څرګندولے شو:

\begin{defn}\ollabel{comparisondef}
  $A$ له $B$ سره \emph{همشمېره} دے، چې $\cardeq{A}{B}$ ئې ليکو،
  هله او يوازې هله چې يوه !!a{bijection} $f \colon A \to B$ وي۔
\end{defn}

\begin{prop}\ollabel{equinumerosityisequi}
همشمېر والی يوه د معادلتوب اړيکه ده۔
\end{prop}

\begin{proof} 
بايد وښيو چې همشمېر والی انعکاسي، متناظر او انتقالي دے۔ فرض
کړئ $A, B$ او $C$ سټونه دي۔

\emph{انعکاسيتوب۔} عيني نقشه $\Id{A} \colon A \to A$، چې د ټولو
$x \in A$ دپاره $\Id{A} (x) = x$ وي، يوه !!a{bijection} ده۔ نو
$\cardeq{A}{A}$۔

\emph{تناظر۔} فرض کړئ $\cardeq{A}{B}$؛ يعنې يوه !!a{bijection}
$f\colon A \to B$ شته۔ ځکه چې $f$ !!{bijective} ده، د هغې معکوس
$f^{-1}$ شته او هغه هم !!{bijective} دے۔ له دې امله
$f^{-1}\colon B \to A$ يوه !!a{bijection} ده، نو $\cardeq{B}{A}$۔

\emph{انتقاليتوب۔} فرض کړئ چې $\cardeq{A}{B}$ او
$\cardeq{B}{C}$؛ يعنې !!{bijection}s $f\colon A \to B$ او
$g\colon B \to C$ شته۔ بيا ترکيب $\comp{f}{g}\colon A \to C$
!!{bijective} دے، نو $\cardeq{A}{C}$۔
\end{proof}

\begin{prop}
که $\cardeq{A}{B}$ وي، نو $A$ هله او يوازې هله !!{enumerable}
دے چې $B$ داسې وي۔
\end{prop}

\begin{editorial}
لاندې ثبوت که \olref[enm]{sec} شامل وي،
\olref[enm]{defn:enumerable} کاروي، او که نۀ وي
\olref[enm-alt]{defn:enumerable} کاروي۔
\end{editorial}

\begin{proof}
فرض کړئ $\cardeq{A}{B}$، نو کومه !!{bijection} $f \colon A
\to B$ شته، او فرض کړئ چې $A$ !!{enumerable} دے۔
\oliflabeldef{sfr:siz:enm:defn:enumerable}{
  بيا يا $A = \emptyset$ دے، يا يوه !!a{surjective} تابع
  $g\colon \PosInt \to A$ شته۔ که $A = \emptyset$ وي، نو
  $B = \emptyset$ هم دے (ګنې يو !!a{element}~$y \in B$ به وے،
  خو هېڅ $x \in A$ به نۀ وے چې $f(x) = y$)۔
  \textbf{د ژباړې د نسخې سمون (OLSIZ-006):} د انګرېزي ثبوت په
  دواړو مشروطو بڼو کښې د تش حالت معادله ناسمې تابعې ته منسوبه
  وه؛ دلته هغه دوه اړخيزې تابعې ته سمه شوې چې د سټونو همشمېر
  والی شاهدي کوي۔ له بلې خوا، که $g\colon \PosInt \to A$
  !!{surjective} وي، نو $\comp{g}{f} \colon \PosInt \to B$
  !!{surjective} ده۔ د دې د ليدلو دپاره $y \in B$ ونيسئ۔ ځکه
  چې $f$ !!{surjective} ده، يو $x \in A$ شته چې $f(x) = y$۔ ځکه
  چې $g$ !!{surjective} ده، يو $n \in \PosInt$ شته چې $g(n) =
  x$۔ له دې امله
\[
(\comp{g}{f})(n) = f(g(n)) = f(x) = y
\]
او نو $\comp{g}{f}$ !!{surjective} ده۔ موږ لرو چې
$\comp{g}{f}$ د~$B$ يوه شمېرنه ده، نو $B$~!!{enumerable} دے۔}
{
  بيا يا $A  = \emptyset$ دے يا يوه !!a{bijection}~$g$ شته چې
  د قيمتونو سټ ئې $A$ او د تعريف ساحه ئې يا $\Nat$ ده يا د طبيعي
  عددونو يو پيلنی تسلسل۔ که $A = \emptyset$ وي، نو
  $B = \emptyset$ هم دے (ګنې به کوم~$y \in B$ وے، خو هېڅ
  $x \in A$ به نۀ وے چې $f(x) = y$)۔ نو فرض کړئ چې زموږ
  !!{bijection}~$g$ شته۔ بيا $\comp{g}{f}$ يوه !!a{bijection}
  ده چې د قيمتونو سټ ئې~$B$ او د تعريف ساحه ئې د~$g$ له ساحې
  سره يوه ده (يعنې يا $\Nat$ يا د هغه يو پيلنی تسلسل)؛ نو $B$
  !!{enumerable} دے۔}

که $B$ !!{enumerable} وي، د~$f$ پر ځاے د !!{bijection}
$f^{-1}\colon B \to A$ په کارولو د استدلال په تکرارولو مومو چې
$A$ !!{enumerable} دے۔
\end{proof}

\begin{prob}
وښيئ چې که $\cardeq{A}{C}$ او $\cardeq{B}{D}$ وي، او
$A \cap B = C \cap D = \emptyset$ وي، نو
$\cardeq{A \cup B}{C \cup D}$۔
\end{prob}

\begin{prob}
وښيئ چې که $A$ نامتناهي او !!{enumerable} وي، نو
$\cardeq{A}{\Nat}$۔
\end{prob}

\end{document}
